\section{Proceso inverso (Reverse Process)}

\subsection{Definición del proceso inverso}

El proceso inverso (Reverse Process) es el núcleo del proceso generativo: parte de ruido puro $x_T \sim \mathcal{N}(0, I)$ y va reduciendo el ruido paso a paso hasta obtener una muestra de datos $x_0$.

La distribución conjunta del proceso inverso se define como:

$$p_\theta(x_0, x_1, \ldots, x_T) = p(x_T) \prod_{t=1}^{T} p_\theta(x_{t-1} | x_t)$$

Donde los parámetros de la distribución de transición inversa en cada paso se parametrizan como una gaussiana:

$$p_\theta(x_{t-1} | x_t) = \mathcal{N}(x_{t-1}; \mu_\theta(x_t, t), \sigma_t^2 I)$$

\begin{importantbox}{Significado físico del proceso inverso}
El proceso inverso es la ``inversión temporal'' del proceso hacia adelante. El proceso hacia adelante añade ruido progresivamente, mientras que el proceso inverso lo elimina paso a paso. La tarea de la red neuronal $\mu_\theta(x_t, t)$ es: dado una imagen con ruido actual $x_t$ y el paso temporal $t$, predecir la media después de un paso de desruido. Obsérvese que la red neuronal necesita conocer la información de $t$, porque los niveles de ruido son diferentes en cada paso temporal.
\end{importantbox}

\subsection{Flujo de generación del proceso inverso}

Dado un modelo entrenado, el proceso de generación es el siguiente:

\begin{enumerate}
    \item Muestrear $x_T \sim \mathcal{N}(0, I)$ desde una distribución gaussiana estándar.
    \item Para $t = T, T-1, \ldots, 1$:
    \begin{itemize}
        \item Calcular $\mu_\theta(x_t, t)$ y $\sigma_t$ usando la red neuronal.
        \item Muestrear $x_{t-1} \sim \mathcal{N}(\mu_\theta(x_t, t), \sigma_t^2 I)$.
    \end{itemize}
    \item Obtener finalmente la muestra generada $x_0$.
\end{enumerate}

\begin{warningbox}{Problema de velocidad de generación}
El proceso inverso requiere ejecutar $T$ pasos de muestreo secuencial (en DDPM, $T = 1000$), y cada paso necesita una propagación hacia adelante de la red neuronal. Esto hace que la velocidad de generación de DDPM sea mucho más lenta que la de GAN (que solo requiere una propagación hacia adelante). Este es también el problema central que trabajos posteriores como DDIM o modelos de consistencia buscan resolver.
\end{warningbox}

\subsection{Resumen de la sección}

El proceso inverso es la única parte del modelo de difusión que necesita aprenderse. Al parametrizar la transición inversa de cada paso como una distribución gaussiana y usar una red neuronal para predecir su media, el modelo aprende a recuperar los datos progresivamente desde el ruido.


